\documentclass[conference]{IEEEtran}
\IEEEoverridecommandlockouts
\usepackage{cite}
\usepackage{amsmath,amssymb,amsfonts}
\usepackage{array}
\usepackage{graphicx}
\usepackage{textcomp}
\usepackage{xcolor}
\usepackage{url}
\newcommand{\stage}[1]{\textsc{#1}}
\begin{document}

\title{Can We Predict the Next Pandemic?\\
Auditable Multistage Forecasting for Surveillance and Preparedness}

\author{\IEEEauthorblockN{Evidence-Grounded Research Proposal}
\IEEEauthorblockA{\textit{Design-only public-health decision-support study}}}

\maketitle

\begin{abstract}
The next pandemic cannot presently be predicted as a single future pathogen, location, and date. That target collapses several distinct transitions: ecological opportunity for spillover, exposure and infection, local amplification, regional dissemination, and the capacity to detect and respond. Yet the failure of a universal oracle does not imply that prediction is futile. This research proposal advances an auditable multistage forecasting framework that estimates calibrated, aggregate region--time risk states for spillover context, local amplification, regional spread readiness, and decision actionability. The framework is designed to help human public-health teams prioritize surveillance and preparedness under declared uncertainty rather than to automate intervention or claim prevention. Its key contribution is an evaluation contract: time-forward and geography-held-out validation, simulated reporting delays, calibration and interval coverage, lead-time and decision-utility analysis, coverage-stratified fairness diagnostics, and an explicit abstention state. The design compares the framework with a pooled risk score and a signal-only baseline while preventing retrospective leakage and random-split overclaiming. Evidence from zoonotic-risk studies, spillover modelling, early-warning reviews, and recent machine-learning critiques supports the component choices and rejects the promise of an exact next-pandemic forecast. No new model has been run and no observed outcomes are reported. The intended deliverable is a responsible research protocol for testing whether stage-specific forecasts can support bounded, accountable preparedness decisions.
\end{abstract}

\begin{IEEEkeywords}
pandemic preparedness, zoonotic spillover, early warning, One Health, calibration, uncertainty, surveillance, decision support
\end{IEEEkeywords}

\section{Introduction}

Pandemic prediction is often posed as a deceptively simple challenge: identify the pathogen, place, and date of the next global emergency. The wording is attractive because it suggests a benchmark with a clear answer. Scientifically, however, it is the wrong benchmark. A future pandemic is not a single event generated by a single mechanism. It is a chain of biological, environmental, institutional, and social transitions. A pathogen may circulate in wildlife without human exposure; a spillover infection may not sustain human-to-human transmission; a local outbreak may be contained; and a regional emergency may fail to become a pandemic because of detection, response, chance, or changing behavior. Conversely, an initially modest signal can be amplified by delayed detection, dense connectivity, inequitable access to care, or fragile logistics.

The correct response is neither to promise a universal prediction engine nor to abandon forecasting. The evidence instead supports a more rigorous objective: estimate and communicate the risk of \emph{specified transitions} for a specified region and future time window, then decide whether that uncertainty justifies a bounded surveillance or preparedness action. Ecological risk mapping, viral and host trait analyses, syndromic surveillance, laboratory confirmation, event-based reports, wastewater signals, mobility context, and health-system information are all potentially useful. They should not be treated as interchangeable evidence of an inevitable pandemic.

This report simulates the current four-agent research process used in the accompanying repository. The Survey Agent constrains claims to documented evidence and a gap ledger. The Idea Agent searches competing scientific directions and selects a falsifiable primary direction. The ExperimentDesign Agent translates that direction into a design-only protocol with variables, baselines, outcome branches, and safety constraints. The Author Agent assembles the source-bound proposal. This is a research plan, not a claim that a new model has predicted an outbreak, prevented a pandemic, or been deployed.

The proposed direction is called \emph{auditable multistage pandemic-escalation forecasting}. It treats the output as a set of calibrated risk states: spillover context, local amplification concern, regional spread readiness, and actionability. A fifth state, abstention, is essential. If surveillance coverage, data rights, outcome definitions, calibration, or equity diagnostics are insufficient, a system should state that it cannot responsibly rank the setting rather than manufacture precision. The aim is assertive but limited: develop an evaluation-ready path to forecasts that can change a named human-reviewed preparedness decision.

\section{Survey Evidence and Research Gap}

\subsection{What the literature permits us to claim}

Zoonotic emergence is not featureless. Olival \emph{et al.} show that host and viral traits can help characterize zoonotic spillover tendencies in mammals \cite{olival2017}. Pandit \emph{et al.} use host and environmental information to predict wildlife reservoirs and vulnerability for zoonotic flaviviruses \cite{pandit2018}. These studies justify prioritization: a surveillance program can decide where characterization effort is likely to be more informative. They do not establish that the next human outbreak can be dated from these covariates.

The broader ecological picture is similarly informative and limited. Allen \emph{et al.} identify global hotspots and correlates of emerging zoonotic diseases, including the role of reporting effort in the observed record \cite{allen2017}. Their discussion is especially important for this proposal: heterogeneous source studies, different reporting practices, and distinct disease ecologies create imprecision. A global model may reveal broad patterns while remaining too coarse for a universal operational prediction. Ecological features should therefore enter as contextual inputs with a coverage and uncertainty record, not as deterministic triggers.

The transition from reservoir to human disease must also be decomposed. Basinski \emph{et al.} combine reservoir ecology and human serosurveys to estimate Lassa virus spillover, illustrating why reservoir presence and human infection cannot be represented by the same label \cite{basinski2021}. Lo Iacono \emph{et al.} formulate zoonotic spillover and subsequent spread as related infection-dynamic processes \cite{loiacono2016}. These works motivate the proposal's first design rule: a high spillover-context estimate, a local-amplification estimate, and a regional-dissemination estimate are different scientific statements, with different evidence needs and intervention implications.

At the detection end, early warning is a system property. Meckawy \emph{et al.} review infectious-disease early-warning systems and show why performance must be considered in relation to data collection, setting, and operational design \cite{meckawy2022}. Event-based reports can complement formal surveillance when they are validated and communicated through appropriate networks \cite{keller2009}. HealthMap illustrates automated classification and visualization of internet media reports for global infectious-disease awareness \cite{freifeld2008}; it is not a substitute for case confirmation. Wastewater-based epidemiology similarly provides a potentially valuable community-level signal but has sampling, interpretation, and implementation limits \cite{polo2020}. A robust framework must preserve source provenance and corroboration status rather than simply add every signal into one opaque score.

The most consequential warning comes from model evaluation. Kawasaki \emph{et al.} identify hidden challenges in evaluating zoonotic-virus spillover-risk machine-learning models, including failure modes driven by sampling structure and biologically structured data \cite{kawasaki2025}. A favourable random train/test split can be an artefact of related observations appearing on both sides of the split. Models should instead be challenged with later periods, held-out geographies, realistic delays, and coverage shifts. Holmes, Rambaut, and Andersen make the complementary strategic argument that resources should not be diverted from surveillance toward inflated prediction promises \cite{holmes2018}. This proposal treats that critique as a specification: forecasting must reinforce verification and preparedness, not replace them.

\begin{table*}[t]
\caption{Evidence-to-claim map used by the Survey Agent.}
\label{tab:evidence}
\centering
\renewcommand{\arraystretch}{1.2}
\begin{tabular}{|p{0.22\textwidth}|p{0.35\textwidth}|p{0.32\textwidth}|}
\hline
\textbf{Evidence line} & \textbf{Permitted inference} & \textbf{Prohibited extrapolation} \\
\hline
Host, virus, and reservoir studies \cite{olival2017,pandit2018,basinski2021} & Prioritize surveillance and distinguish reservoir, exposure, and infection layers. & Predict a specific future pathogen, date, or spillover event. \\
\hline
Ecological emergence analyses \cite{allen2017} & Use land-use, biodiversity, climate, and reporting context as uncertainty-qualified inputs. & Treat reported global association as a universal causal rule. \\
\hline
Early-warning and event surveillance \cite{meckawy2022,keller2009,freifeld2008,polo2020} & Combine and corroborate signals in a human-reviewed surveillance system. & Equate an anomaly with a confirmed outbreak or prevention outcome. \\
\hline
ML evaluation critiques \cite{kawasaki2025} & Require temporal, geographic, and bias-aware validation with calibration. & Use random-split accuracy as proof of future operational value. \\
\hline
Surveillance-first argument \cite{holmes2018} & Connect forecasts to verification, preparedness capacity, and accountable decisions. & Substitute a prediction narrative for surveillance investment. \\
\hline
\end{tabular}
\end{table*}

\subsection{Gap ledger}

Five gaps organize the research question. \stage{Gap-01} is target conflation: the field and public discourse often call several distinct products ``pandemic prediction.'' \stage{Gap-02} is stage attribution: even when a risk score is high, it may be unclear whether the concern comes from ecological exposure, local transmission, mobility, weak detection, or a combination. \stage{Gap-03} is evaluation realism: models are frequently under-tested against future time periods, new locations, reporting lags, and label revisions. \stage{Gap-04} is decision linkage: a rank without a threshold, action owner, lead-time need, and false-alert cost is not an operational product. \stage{Gap-05} is equity and observability: thin surveillance can create low apparent risk while concealing a lack of evidence.

These gaps are connected. A pooled score can look accurate because it learns reporting effort, while its users assume it estimates biological hazard. Its apparent accuracy can then be used to direct attention away from poorly observed settings. Stage decomposition does not erase this risk; it makes the risk inspectable and provides a principled path to abstention.

\section{Idea Agent: Direction Search and Selection}

The Idea Agent explored four routes. The first, reservoir and viral discovery prioritization, is scientifically valuable but stops before an outbreak-readiness decision. The second, real-time event-based anomaly detection, can shorten detection time but is weak on upstream conditions and can create false alerts without corroboration. The third, a universal artificial-intelligence predictor of the next pandemic, was rejected. It fails the survey constraints because it collapses targets, has no credible universal label, offers weak falsification, and invites unjustified automation.

The selected route is the fourth: an auditable multistage framework. It is built through four conceptual edits. First, an ecological-hotspot seed is retained as a contextual input rather than as a final prediction. Second, a single fused score is rejected because a user cannot tell whether it responds to a changing host environment, an incomplete clinical feed, or a retrospective reporting revision. Third, the target is decomposed into stage-specific states. Fourth, the states are bound to an action table with uncertainty, data freshness, equity flags, and abstention. The selected idea preserves the useful information in upstream and real-time signals while preventing their semantic collapse.

\begin{figure*}[t]
\centering
\fbox{\begin{minipage}{0.93\textwidth}
\centering\small
\textbf{Survey Agent} \quad $\longrightarrow$ \quad \textbf{Idea Agent} \quad $\longrightarrow$ \quad \textbf{ExperimentDesign Agent} \quad $\longrightarrow$ \quad \textbf{Author Agent} \\
\vspace{0.35em}
Evidence map, source limits, and gap ledger \quad $\longrightarrow$ \quad competing directions, debate, selected D2 \quad $\longrightarrow$ \quad variables, validation, safety, decision branches \quad $\longrightarrow$ \quad source-bound IEEE proposal \\
\vspace{0.85em}
\fbox{\parbox{0.83\textwidth}{\centering\textbf{Selected scientific target:} calibrated aggregate risk states for \emph{spillover context} $\rightarrow$ \emph{local amplification} $\rightarrow$ \emph{regional spread readiness} $\rightarrow$ \emph{human-reviewed actionability}. At every transition: data-quality record, uncertainty interval, equity diagnostic, and \textbf{abstain} option.}}
\end{minipage}}
\caption{Traceable four-agent workflow and the selected multistage prediction target. The workflow produces a proposal; it reports no new observed results.}
\label{fig:workflow}
\end{figure*}

The central hypothesis is deliberately demanding: for predeclared decisions, a stage-specific framework will be more interpretable and better calibrated than a strong pooled-score baseline. This is not a claim that decomposition must win. It may fail if the stages are poorly observed, if labels cannot be defined consistently, if errors are too correlated, or if the extra structure reduces lead-time utility. Those outcomes are informative and must be reported rather than repaired through post hoc target changes.

\section{Research Questions and Formal Hypotheses}

The proposal poses four research questions.

\textbf{RQ1:} Can stage-specific inputs and outputs make the mechanism of a high-risk assessment auditable? \textbf{H1} predicts that reviewers will be able to trace a sampled output to a stated stage, source age, uncertainty interval, and decision rule more consistently than for a pooled score.

\textbf{RQ2:} Does stage-aware modelling preserve calibration under realistic deployment shifts? \textbf{H2} predicts better calibration and interval coverage on time-forward and geography-held-out evaluations than a pooled model evaluated under the same information constraints. The criterion is not raw area-under-the-curve alone; a highly discriminative but overconfident system can make worse decisions than a calibrated one.

\textbf{RQ3:} Does a risk state improve a decision with a declared lead-time requirement and false-alert cost? \textbf{H3} predicts higher net benefit or lower expected loss when a human reviewer uses the stage output plus an actionability state, compared with a threshold applied to the pooled score or a signal-only baseline.

\textbf{RQ4:} When should the system not rank a setting? \textbf{H4} predicts that explicit reporting-delay, coverage, and research-effort diagnostics will identify cases in which abstention is more responsible than a false-precision score. A high abstention rate is not automatically a failure: it can reveal a surveillance investment need.

Let region $r$ and decision time $t$ define the unit of analysis. Let $k \in \{1,2,3\}$ denote spillover context, local amplification, and regional spread readiness. The proposed framework estimates
\begin{equation}
p^{(k)}_{r,t,h}=\Pr\left(Y^{(k)}_{r,t+h}=1\mid X^{(k)}_{r,\leq t},Q_{r,t}\right),
\label{eq:stageprob}
\end{equation}
where $h$ is a predeclared decision horizon, $X^{(k)}$ contains only information available at or before $t$, and $Q$ records quality, timeliness, coverage, and revision indicators. Crucially, the equation does not imply that the $Y^{(k)}$ labels are equally observable or causally complete. It defines a testable probability target conditional on explicit label and data-governance choices.

The actionability output is a decision function rather than a biological probability:
\begin{equation}
a_{r,t}=g\left(p^{(1)}_{r,t,h},p^{(2)}_{r,t,h},p^{(3)}_{r,t,h},Q_{r,t},C,L\right),
\label{eq:actionability}
\end{equation}
where $C$ is a documented cost structure and $L$ is a required lead time. The function must return an abstention option whenever data, calibration, governance, or fairness requirements are not met. Neither Eq.~\eqref{eq:stageprob} nor Eq.~\eqref{eq:actionability} is an observed result; both are a preregistrable design object.

\section{ExperimentDesign Agent: Study Design}

\subsection{Execution boundary and units}

This is a \textbf{design-only} computational and public-health decision-support protocol. It does not collect new specimens, manipulate pathogens, enrol participants, perform animal work, access individual movement traces, or trigger a field intervention. If later implemented, all nonpublic data access, jurisdictional partnerships, and prospective use require appropriate ethics, legal, data-governance, and community review.

The unit is an aggregate region-by-predeclared-time-window. Spatial granularity must be chosen to avoid re-identification and stigmatizing precision. The time horizon must be selected before model fitting to match a specific action, such as reviewing surveillance capacity, pre-positioning nonclinical supplies, checking laboratory throughput, or increasing corroborating surveillance. The protocol does not authorize intrusive or punitive measures.

Inputs are represented as versioned data cards, not an indiscriminate feature list. Spillover-context cards may summarize authorized host/virus surveillance, land-use, climate anomaly, and aggregate exposure context. Local-amplification cards may include validated aggregate laboratory, syndromic, or event indicators. Spread-readiness cards may describe aggregate connectivity and health-system surge context. Every card records provenance, access rights, spatial and temporal coverage, refresh date, revision behavior, known bias, and intended interpretation. Data quality itself is modelled as an output condition; missingness cannot silently become apparent safety.

\subsection{Comparators and outcome definitions}

The primary comparator is a calibrated pooled score that uses the same eligible aggregate information but emits one risk estimate. This is intentionally a serious baseline; the multistage proposal should not win merely because it is compared with a weak heuristic. The second comparator is a signal-only baseline based on contemporary clinical, laboratory, or event anomaly evidence. It tests whether the upstream and decision layers add value beyond detection alone. An oracle comparator is explicitly excluded because it would use future confirmations, retrospectively corrected labels, or information unavailable at the decision time.

Outcome definitions will be negotiated with public-health partners before implementation. Examples include predeclared, aggregate stage labels for evidence of an unusual zoonotic exposure context, a locally sustained transmission concern, or a regional dissemination-relevant episode. These examples are not labels for individual persons or declarations of a future pandemic. Definition changes after outcome inspection must be versioned and analyzed as a new protocol, not silently merged with the original study.

\begin{table*}[t]
\caption{Design variables, safeguards, and the reason each is retained.}
\label{tab:variables}
\centering
\renewcommand{\arraystretch}{1.15}
\begin{tabular}{|p{0.18\textwidth}|p{0.29\textwidth}|p{0.25\textwidth}|p{0.17\textwidth}|}
\hline
\textbf{Construct} & \textbf{Authorized aggregate representation} & \textbf{Output or evaluation role} & \textbf{Guardrail} \\
\hline
Spillover context & Host/virus-surveillance summaries, ecological change, climate anomaly, exposure proxies & $p^{(1)}$ and uncertainty interval & No wildlife targeting or individual exposure score. \\
\hline
Local amplification & Validated aggregate laboratory, syndromic, and corroborated event indicators & $p^{(2)}$, source freshness, corroboration state & Anomaly is not a confirmed case count. \\
\hline
Regional readiness & Aggregate connectivity, reporting latency, surge-capacity context & $p^{(3)}$ and preparedness relevance & No individual mobility tracking. \\
\hline
Data quality & Coverage, revision, missingness, effort, source age, access constraint & reliability card and abstention trigger & Never impute away absence of surveillance. \\
\hline
Decision utility & Named action owner, cost, lead-time need, opportunity cost & actionability state $a_{r,t}$ & No automatic operational action. \\
\hline
\end{tabular}
\end{table*}

\subsection{Model and communication architecture}

The study will compare implementation families rather than prespecify a fashionable algorithm. Each stage may use a transparent statistical model, a constrained machine-learning model, or an ensemble if the justification, fitting information set, and calibration method are documented. The same feature time cut-off applies to every model. If models are linked, the interface must expose rather than hide uncertainty propagation. For example, a low-confidence upstream stage cannot be converted into a confident downstream action solely through model complexity.

The output is a compact risk card. It includes the region and horizon; stage estimates and intervals; data freshness and coverage; the strongest corroborated evidence class; calibration status for comparable settings; an action class; and a plain-language reason for abstention where applicable. The card is an input to a human decision meeting, not an alert to the public. No public-facing map is generated by this design until its harms, interpretability, and communication governance are separately reviewed.

\section{Evaluation, Decision Logic, and Failure Analysis}

\subsection{Deployment-realistic validation}

Validation begins with a protocol freeze: data dictionary, source eligibility, regional resolution, stage outcomes, missing-data treatment, decision horizons, actions, thresholds, and analysis code plan. Historical data are then arranged in sequential training and later evaluation windows. Entire geographic units are held out to test transportability. Data releases are reconstructed with realistic reporting lags and revisions. This prevents future information from entering past decisions and exposes the operational effect of stale or incomplete streams.

Primary measures are calibration slope and intercept, proper scoring rules, interval coverage, decision-curve net benefit, lead-time distribution, abstention rate, and error disparity across coverage strata. Discrimination may be reported but is secondary. A model that ranks correctly while assigning untrustworthy probabilities cannot support threshold decisions. Metrics are reported per stage, per horizon, per data-quality stratum, and per geography where reporting is ethically and statistically justified. Aggregate summaries must not conceal a subgroup with systematically worse calibration.

The protocol includes a reviewer-facing attribution audit. For a stratified sample of high-risk, low-risk, and abstained outputs, reviewers check whether the named stage follows from the available source card, whether input dates meet the horizon rule, whether the interval matches the calibration record, and whether the suggested action is permissible under the action table. This audit tests the proposal's core promise: a decision maker should be able to understand why the system expressed concern and why it may refuse to act.

\begin{figure}[t]
\centering
\fbox{\begin{minipage}{0.94\columnwidth}
\centering\footnotesize
\textbf{1. Spillover context}\\
Ecological / host / virus evidence\\
$\Downarrow$ (uncertainty and coverage)\\
\textbf{2. Local amplification}\\
Laboratory / syndromic / event corroboration\\
$\Downarrow$ (calibration and latency)\\
\textbf{3. Regional spread readiness}\\
Connectivity / capacity / reporting context\\
$\Downarrow$ (utility, cost, and lead time)\\
\textbf{4. Human-reviewed actionability}\\
\fbox{prepare} \quad \fbox{corroborate} \quad \fbox{monitor} \quad \fbox{abstain}
\end{minipage}}
\caption{Decision path. Every transition preserves data-quality and uncertainty information; the final class never triggers autonomous action.}
\label{fig:decisionpath}
\end{figure}

\subsection{Conditional conclusions}

The study has four permitted conclusion labels. \emph{Stage-specific signal supported} requires stable held-out calibration and a successful traceability audit; it means that the stage representation warrants further monitored evaluation. \emph{Actionable lead time supported} additionally requires superior utility at a predeclared false-alert cost and a decision-relevant lead time; it means that a limited, human-reviewed preparedness pilot may be considered. \emph{Transportability insufficient} is required when time or geography transfer fails materially, even if random-split performance looked favorable. \emph{Definition or governance abstain} is required when labels, data rights, equity diagnostics, or accountable action rules are inadequate.

Several negative findings would be valuable. The multistage system may not improve calibration over a well-calibrated pooled baseline. A component may add apparent discrimination but reduce decision utility because its reporting latency is too large. Geographic shifts may expose hidden reliance on local surveillance practice. Equity diagnostics may reveal that the model is most confident where data are already abundant. Each outcome identifies a concrete response: simplify the architecture, collect or govern better data, localize the model, adjust the decision horizon, or abstain. None should be reframed as proof that a pandemic was unpredictable in principle.

\section{Prospective Shadow Evaluation and Audit Protocol}

Retrospective analysis is necessary but not sufficient. Data streams are revised, reporting incentives change, and apparent labels may be cleaner in a historical archive than they were at the time a decision had to be made. Before any claim of operational readiness, the design therefore requires a prospective shadow phase. During this phase, authorized aggregate inputs are frozen on a regular schedule and transformed into risk cards without changing public-health operations. The cards are available only to an approved review group under a protocol that separates model evaluation from real-time response. Later outcome information is used to assess calibration, lead time, and abstention decisions, but never to rewrite the card that was originally issued.

The shadow phase is deliberately not a passive dashboard exercise. Each evaluation cycle records five auditable items: the data version available at the cut-off; the stage estimates and intervals; the reliability and coverage flags; the action class or abstention; and the reviewer rationale. The reviewer may judge that a card is internally inconsistent, inadequately explained, or governed by an unacceptable threshold. Such judgments are retained as structured audit data rather than treated as inconvenient anecdote. This gives the study a way to distinguish a statistical failure from an interface, governance, or decision-specification failure.

An action class has meaning only when a responsible party, required lead time, and opportunity cost are specified. For example, a \emph{corroborate} class may authorize a human team to verify whether authorized aggregate sources agree and whether routine laboratory or sentinel capacity is sufficient. A \emph{prepare} class may initiate a nonclinical capacity review or logistics check already permitted by the jurisdiction. A \emph{monitor} class can preserve routine surveillance while recording why escalation is not justified. The \emph{abstain} class states that the system cannot responsibly recommend any of the preceding options. It is not an instruction to ignore risk; it is an instruction to resolve a data, definition, authority, or equity deficit before using a quantitative rank.

The action table is frozen before the shadow period, including expected benefits, false-alert costs, and prohibited uses. If different jurisdictions have different capacities, the framework must not hide that difference behind a common numerical threshold. It may use locally approved action tables while maintaining a common reporting structure: source freshness, uncertainty, calibration evidence, coverage diagnostic, and named reviewer. This is essential because a technically identical risk estimate can imply different responsible decisions in systems with different laboratory throughput, health-care access, and governance authority.

\begin{table}[t]
\caption{Predeclared shadow-evaluation audit record.}
\label{tab:audit}
\centering
\renewcommand{\arraystretch}{1.13}
\begin{tabular}{|p{0.26\columnwidth}|p{0.62\columnwidth}|}
\hline
\textbf{Record field} & \textbf{Required content} \\
\hline
Information cut-off & Time stamp, source versions, late or revised feeds excluded from the decision. \\
\hline
Stage card & Estimates, uncertainty intervals, source-quality flags, and coverage stratum. \\
\hline
Decision mapping & Prepared action class, required lead time, cost assumption, and named human owner. \\
\hline
Audit rationale & Whether stage attribution, uncertainty, and action explanation are internally consistent. \\
\hline
Later evaluation & Matched predeclared outcome, calibration contribution, utility contribution, and any protocol deviation. \\
\hline
\end{tabular}
\end{table}

The protocol also defines a model-change discipline. Updates may be triggered by prespecified drift criteria, data-source retirement, a calibration failure, or a documented change in outcome definition. An update cannot be justified simply because it makes the model appear better on a recently inspected outcome set. Each version has its own training information boundary, validation report, source cards, decision table, and fairness diagnostics. A new version is evaluated first in retrospective time-forward tests and then, where approved, in a new shadow period. This preserves a chain of evidence for claims about improvement.

Finally, the shadow phase includes a red-team style failure review conducted by the participating human experts. The review asks whether a confident card relied on stale sources, whether a high score was actually a proxy for surveillance intensity, whether a low score was misread as safety, whether a stage was named without adequate evidence, and whether a recommended action exceeded the authority or capacity of the proposed user. The review has no mandate to discover operational vulnerabilities or to direct interventions. Its role is to test whether the system's assurances of traceability and restraint survive realistic decision pressure.

\section{Evaluation Scenarios and Interpretation Rules}

The protocol evaluates a framework by the quality of its decisions under distinct information conditions, not by a single average score. Scenario analysis is therefore part of the primary design. The scenarios are synthetic evaluation states constructed from authorized aggregate historical records and documented perturbations such as a delayed report, a missing data feed, a changed coverage pattern, or a disagreement between signal classes. They do not simulate pathogens, generate biological data, or issue live alerts. Their purpose is to ensure that the same numerical output is not interpreted identically when its evidentiary basis changes.

The first scenario is \emph{upstream concern with weak downstream corroboration}. Ecological, host, or viral context may indicate elevated spillover relevance while clinical, laboratory, and event signals remain ordinary or unavailable. A pooled model could convert this mixture into an alarming global rank. The multistage model must instead report a bounded spillover-context state, low or uncertain amplification evidence, the age and coverage of downstream sources, and a decision class no stronger than a human-reviewed corroboration or monitoring action. A ``prepare'' action is allowed only if the action table shows that it is proportional to the stated probability, capacity context, and false-alert cost. This scenario tests whether the framework preserves the distinction between prioritizing surveillance and declaring an outbreak.

The second scenario is \emph{downstream anomaly with uncertain upstream mechanism}. A validated aggregate clinical, laboratory, or event anomaly may deserve rapid corroboration even when a reservoir or ecological explanation is unknown. The design must not suppress such signals because an upstream score is low. It should report the local-amplification state separately, flag the absence of an established spillover mechanism, and make clear that the output concerns contemporaneous detection rather than a prior prediction. This scenario tests whether stage decomposition makes a framework more responsive rather than more bureaucratic. It also prevents a retrospective narrative in which every detected anomaly is incorrectly attributed to an ecological predictor after the fact.

The third scenario is \emph{apparent low risk under sparse observation}. A region may have weak laboratory coverage, irregular reporting, limited historical research effort, or delayed data releases. A naïve model can interpret few reports as reassurance. In the proposed framework, the quality state dominates: low coverage widens intervals, invalidates comparable calibration evidence, and can force the abstain class. This is a scientifically strong output because it explicitly identifies a limit of inference. The appropriate response is not automatic escalation against a community; it is a human governance decision about whether and how surveillance capacity should be supported.

\begin{table*}[t]
\caption{Predeclared interpretation rules for design-only evaluation scenarios.}
\label{tab:scenarios}
\centering
\renewcommand{\arraystretch}{1.15}
\begin{tabular}{|p{0.21\textwidth}|p{0.27\textwidth}|p{0.22\textwidth}|p{0.17\textwidth}|}
\hline
\textbf{Scenario} & \textbf{Expected stage representation} & \textbf{Evaluation question} & \textbf{Maximum permissible response class} \\
\hline
Upstream concern, weak downstream corroboration & Elevated spillover context; uncertain amplification and readiness; source-age card visible & Does the system avoid converting priority context into an outbreak claim? & Monitor or human-reviewed corroboration unless the predeclared utility rule supports more. \\
\hline
Downstream anomaly, uncertain upstream mechanism & Amplification concern shown independently; unknown mechanism retained & Does the system detect a relevant anomaly without inventing a spillover explanation? & Human-reviewed corroboration and capacity review. \\
\hline
Sparse observation, apparent low risk & Wide intervals; poor coverage; calibration non-comparability; abstention & Does missing surveillance lower confidence rather than lower risk by default? & Abstain from a quantitative recommendation. \\
\hline
Conflicting signals with late revision & Competing source cards; revision flag; revised uncertainty and actionability & Does the card remain stable about what was known at decision time? & Monitor, corroborate, or abstain according to the frozen version. \\
\hline
\end{tabular}
\end{table*}

Interpretation rules also protect against hindsight bias. A later-confirmed event cannot make an earlier probability calibrated by definition. Each card is scored against an outcome chosen in advance, using the record available at issuance. Reviewers may compare the original and revised card to quantify the impact of information latency, but they must not replace the original decision-time record. Similarly, a model cannot receive credit for a broad warning if its proposed action class was vague enough to be unfalsifiable. The action must be sufficiently specified that a later evaluator can decide whether the lead time and net benefit criteria were met.

The scenarios provide a direct test of the selected Idea Agent claim. If a pooled score communicates the same information, is as well calibrated, yields equal or better decision utility, and has fewer coverage-related harms, the multistage architecture adds unjustified complexity and should not be adopted. If the multistage system makes uncertainty and mechanism more legible, identifies conditions for abstention, and supports better bounded decisions under declared costs, then it earns its additional structure. This is a demanding standard, but it is appropriate for a field where a false sense of prediction can be as harmful as a missed signal.

\section{Governance, Equity, and Limitations}

Forecasting infectious-disease risk has direct social consequences. A regional score can stigmatize communities, influence travel and investment decisions, or create misplaced confidence. The technical protocol must therefore include a policy and governance protocol. Each participating jurisdiction retains authority over data sharing, interpretation, and response. The intended use, prohibited use, escalation pathway, appeal mechanism, and communication strategy are documented before any prospective shadow deployment. Outputs are aggregate and human-reviewed; the design excludes individual risk scoring and autonomous response.

Equity is not a post-processing metric. Surveillance absence may be driven by historical underinvestment, conflict, barriers to testing, or data sovereignty constraints. Treating an absence of reports as evidence of low risk can deepen the very blind spots that make prediction harder. The framework therefore displays coverage and reporting-delay diagnostics alongside biological or epidemiological estimates. A model that abstains in thinly observed settings is not equitable by itself, but it is more honest than a rank that masks lack of observability. Any eventual deployment must pair abstention analysis with investment and governance decisions made by affected institutions and communities.

There are technical limitations. Stage labels may be uncertain or nonidentical across pathogens and jurisdictions. Causal claims cannot be inferred merely because a covariate improves a forecast. Data-sharing rules may preclude important inputs. Probabilistic calibration can degrade under political, behavioral, climatic, or technological shocks. The proposed architecture does not resolve those limitations; it makes them explicit and testable. The framework should be abandoned or restricted if it cannot maintain calibration, transparency, fairness diagnostics, and actionability under the conditions in which it would be used.

The proposal also has a scope limitation: it does not settle philosophical questions about what counts as a pandemic, nor does it prioritize one disease family over another. It instead establishes a general decision-support pattern. This pattern can be adapted only after subject-matter experts define the biological target, authorized inputs, and permissible actions for a particular setting. That restraint is not passivity. It is what makes a strong predictive claim falsifiable and accountable.

\section{Conclusion}

The next pandemic should not be presented as a riddle with one hidden answer. Exact pathogen, time, and place prediction is not supported by the evidence reviewed here, and a single black-box score would make its own failure mechanisms invisible. At the same time, a surveillance-first position need not reject forecasting. The scientifically defensible opportunity is to forecast limited, stage-specific risk states; to calibrate them under the temporal, geographic, and reporting shifts that deployment will face; and to connect them to a named human-reviewed preparedness action or an explicit abstention.

This four-agent proposal makes that opportunity concrete. Its Survey Agent sets evidence limits; its Idea Agent rejects the universal-oracle narrative and selects an auditable alternative; its ExperimentDesign Agent supplies a design-only, bias-aware evaluation protocol; and its Author Agent preserves the distinction between expected and observed outcomes. The decisive test is practical: can the framework give a justified, timely, and equitable reason to corroborate, prepare, monitor, or abstain? If it cannot, it has not predicted what public health needs. If it can, it offers a meaningful path from uncertainty to preparedness without pretending that uncertainty has disappeared.

\begin{thebibliography}{00}
\bibitem{olival2017} K. J. Olival \emph{et al.}, ``Host and viral traits predict zoonotic spillover from mammals,'' \emph{Nature}, vol. 546, pp. 646--650, 2017, doi: 10.1038/nature22975.
\bibitem{pandit2018} P. S. Pandit \emph{et al.}, ``Predicting wildlife reservoirs and global vulnerability to zoonotic flaviviruses,'' \emph{Nature Communications}, vol. 9, art. 5425, 2018, doi: 10.1038/s41467-018-07896-2.
\bibitem{allen2017} T. Allen \emph{et al.}, ``Global hotspots and correlates of emerging zoonotic diseases,'' \emph{Nature Communications}, vol. 8, art. 1124, 2017, doi: 10.1038/s41467-017-00923-8.
\bibitem{basinski2021} A. J. Basinski \emph{et al.}, ``Bridging the gap: Using reservoir ecology and human serosurveys to estimate Lassa virus spillover in West Africa,'' \emph{PLoS Computational Biology}, vol. 17, no. 3, e1008811, 2021, doi: 10.1371/journal.pcbi.1008811.
\bibitem{loiacono2016} G. Lo Iacono \emph{et al.}, ``A unified framework for the infection dynamics of zoonotic spillover and spread,'' \emph{PLoS Neglected Tropical Diseases}, vol. 10, no. 9, e0004957, 2016, doi: 10.1371/journal.pntd.0004957.
\bibitem{meckawy2022} M. A. Meckawy \emph{et al.}, ``Effectiveness of early warning systems in the detection of infectious diseases outbreaks: a systematic review,'' \emph{BMC Public Health}, vol. 22, art. 2210, 2022, doi: 10.1186/s12889-022-14625-4.
\bibitem{keller2009} M. Keller \emph{et al.}, ``Use of unstructured event-based reports for global infectious disease surveillance,'' \emph{Emerging Infectious Diseases}, vol. 15, no. 5, pp. 689--695, 2009, doi: 10.3201/eid1505.081114.
\bibitem{freifeld2008} C. C. Freifeld \emph{et al.}, ``HealthMap: global infectious disease monitoring through automated classification and visualization of Internet media reports,'' \emph{Journal of the American Medical Informatics Association}, vol. 15, no. 2, pp. 150--157, 2008, doi: 10.1197/jamia.M2544.
\bibitem{polo2020} D. Polo \emph{et al.}, ``Making waves: Wastewater-based epidemiology for COVID-19---approaches and challenges for surveillance and prediction,'' \emph{Water Research}, vol. 186, art. 116404, 2020, doi: 10.1016/j.watres.2020.116404.
\bibitem{kawasaki2025} K. Kawasaki, Y. Suzuki, and K. Hamada, ``Hidden challenges in evaluating spillover risk of zoonotic viruses using machine learning models,'' \emph{Communications Medicine}, vol. 5, art. 251, 2025, doi: 10.1038/s43856-025-00903-w.
\bibitem{holmes2018} E. C. Holmes, A. Rambaut, and K. G. Andersen, ``Pandemics: spend on surveillance, not prediction,'' \emph{Nature}, vol. 558, pp. 180--182, 2018, doi: 10.1038/d41586-018-05373-w.
\end{thebibliography}

\end{document}
